\section{G}

% geben (to give)
\begin{verbentry}{
  style = {dat},
  toc = {geben (to give)},
  name = {geben},
  english = {to give},
  type = {irregular, + Dativ},
  auxiliary = {haben}
}
 
    
     \vspace{5pt}

    \hrule
    
    \vspace{5pt}

  \begin{tabular}{rl}
     \multicolumn{2}{l}{\textbf{PRÄSENS} }\\
    ich & \textbf{gebe}\\
    du & \textbf{gibst}\\
    er/sie/es & \textbf{gibt}\\
    wir & \textbf{geben}\\
    ihr & \textbf{gebt}\\
    sie/Sie & \textbf{geben}\\
  \end{tabular}
  \hfill
     \begin{tabular}{rl}
    \multicolumn{2}{l}{\textbf{PRÄTERITUM} }\\
    ich & \textbf{gab}\\
    du & \textbf{gabst}\\
    er/sie/es & \textbf{gab}\\
    wir & \textbf{gaben}\\
    ihr & \textbf{gabt}\\
    sie/Sie & \textbf{gaben}\\
  \end{tabular}
  \hfill
  \begin{tabular}{rl}
     \multicolumn{2}{l}{\textbf{IMPERATIV} }\\
   & \textbf{gib (du)!}\\
   & \textbf{geben wir!}\\
   & \textbf{gebt (ihr)!}\\
   & \textbf{geben Sie!}\\
   & \vspace{10pt}
  \end{tabular}

    \vspace{10pt}

 \begin{tabular}{rl}
     \multicolumn{2}{l}{\textbf{PERFEKT}}\\
   & haben + \textbf{gegeben}\\
  \end{tabular}
\hfill
   \begin{tabular}{rl}
   
    \multicolumn{2}{l}{\textbf{FUTURE I} }\\
  &  werden + \textbf{geben}\\
 
  \end{tabular}
\hfill
   \begin{tabular}{rl}
      \multicolumn{2}{l}{\textbf{PLUSQUAMPERFEKT}}\\
  &  hatten + \textbf{gegeben}\\
  \end{tabular}
  
  
\vspace{10pt}

\begin{tabular}{lll}
\multicolumn{3}{l}{\textbf{MIT PRÄPOSITIONEN}}\\
& \textbf{---} & \\
\end{tabular}
\hfill
\begin{tabular}{lll}
\multicolumn{3}{l}{\textbf{TRENNBARE VERBEN}}\\
& \textbf{---} & \\
\end{tabular}
\vspace{-5pt}
\end{verbentry}

% gefallen (to like)
\begin{verbentry}{
  style = {dat},
  toc = {gefallen (to like)},
  name = {gefallen},
  english = {to like},
  type = {irregular, + Dativ},
  auxiliary = {haben}
}
 
    
     \vspace{5pt}

    \hrule
    
    \vspace{5pt}

  \begin{tabular}{rl}
     \multicolumn{2}{l}{\textbf{PRÄSENS} }\\
    ich & \textbf{gefalle}\\
    du & \textbf{gefällst}\\
    er/sie/es & \textbf{gefällt}\\
    wir & \textbf{gefallen}\\
    ihr & \textbf{gefallt}\\
    sie/Sie & \textbf{gefallen}\\
  \end{tabular}
  \hfill
     \begin{tabular}{rl}
    \multicolumn{2}{l}{\textbf{PRÄTERITUM} }\\
    ich & \textbf{gefiel}\\
    du & \textbf{gefielst}\\
    er/sie/es & \textbf{gefiel}\\
    wir & \textbf{gefielen}\\
    ihr & \textbf{gefielt}\\
    sie/Sie & \textbf{gefielen}\\
  \end{tabular}
  \hfill
  \begin{tabular}{rl}
     \multicolumn{2}{l}{\textbf{IMPERATIV} }\\
   & \textbf{gefall(e) (du)!}\\
   & \textbf{gefallen wir!}\\
   & \textbf{gefallt (ihr)!}\\
   & \textbf{gefallen Sie!}\\
   & \vspace{10pt}
  \end{tabular}

    \vspace{10pt}

 \begin{tabular}{rl}
     \multicolumn{2}{l}{\textbf{PERFEKT}}\\
  &  haben + \textbf{gefallen}\\
  \end{tabular}
\hfill
   \begin{tabular}{rl}
   
    \multicolumn{2}{l}{\textbf{FUTURE I} }\\
  &  werden + \textbf{gefallen}\\
 
  \end{tabular}
\hfill
   \begin{tabular}{rl}
      \multicolumn{2}{l}{\textbf{PLUSQUAMPERFEKT}}\\
   & hatten + \textbf{gefallen}\\
  \end{tabular}
  
  
\vspace{10pt}

\begin{tabular}{lll}
\multicolumn{3}{l}{\textbf{MIT PRÄPOSITIONEN}}\\
& \textbf{---} & \\
\end{tabular}
\hfill
\begin{tabular}{lll}
\multicolumn{3}{l}{\textbf{TRENNBARE VERBEN}}\\
& \textbf{---} & \\
\end{tabular}
\vspace{-5pt}
\end{verbentry}

% gehören (to belong)
\begin{verbentry}{
  style = {dat},
  toc = {gehören (to belong)},
  name = {gehören},
  english = {to belong},
  type = {regular, + Dativ},
  auxiliary = {haben}
}
 
    
     \vspace{5pt}

    \hrule
    
    \vspace{5pt}

  \begin{tabular}{rl}
     \multicolumn{2}{l}{\textbf{PRÄSENS} }\\
    ich & \textbf{gehöre}\\
    du & \textbf{gehörst}\\
    er/sie/es & \textbf{gehört}\\
    wir & \textbf{gehören}\\
    ihr & \textbf{gehört}\\
    sie/Sie & \textbf{gehören}\\
  \end{tabular}
  \hfill
     \begin{tabular}{rl}
    \multicolumn{2}{l}{\textbf{PRÄTERITUM} }\\
    ich & \textbf{gehörte}\\
    du & \textbf{gehörtest}\\
    er/sie/es & \textbf{gehörte}\\
    wir & \textbf{gehörten}\\
    ihr & \textbf{gehörtet}\\
    sie/Sie & \textbf{gehörten}\\
  \end{tabular}
  \hfill
  \begin{tabular}{rl}
     \multicolumn{2}{l}{\textbf{IMPERATIV} }\\
   & \textbf{gehör(e) (du)!}\\
   & \textbf{gehören wir!}\\
   & \textbf{gehört (ihr)!}\\
   & \textbf{gehören Sie!}\\
   & \vspace{10pt}
  \end{tabular}

    \vspace{10pt}

 \begin{tabular}{rl}
     \multicolumn{2}{l}{\textbf{PERFEKT}}\\
  &  haben + \textbf{gehört}\\
  \end{tabular}
\hfill
   \begin{tabular}{rl}
   
    \multicolumn{2}{l}{\textbf{FUTURE I} }\\
   & werden + \textbf{gehören}\\
 
  \end{tabular}
\hfill
   \begin{tabular}{rl}
      \multicolumn{2}{l}{\textbf{PLUSQUAMPERFEKT}}\\
   & hatten + \textbf{gehört}\\
  \end{tabular}
  
  
\vspace{10pt}

\begin{tabular}{lll}
\multicolumn{3}{l}{\textbf{MIT PRÄPOSITIONEN}}\\
& \textbf{---} & \\
\end{tabular}
\hfill
\begin{tabular}{lll}
\multicolumn{3}{l}{\textbf{TRENNBARE VERBEN}}\\
& \textbf{---} & \\
\end{tabular}
\vspace{-5pt}
\end{verbentry}

% gehen (to go)
\begin{verbentry}{
  style = {acc},
  toc = {gehen (to go)},
  name = {gehen},
  english = {to go},
  type = {irregular},
  auxiliary = {sein}
}
 
    
     \vspace{5pt}

    \hrule
    
    \vspace{5pt}

  \begin{tabular}{rl}
     \multicolumn{2}{l}{\textbf{PRÄSENS} }\\
    ich & \textbf{gehe}\\
    du & \textbf{gehst}\\
    er/sie/es & \textbf{geht}\\
    wir & \textbf{gehen}\\
    ihr & \textbf{geht}\\
    sie/Sie & \textbf{gehen}\\
  \end{tabular}
  \hfill
     \begin{tabular}{rl}
    \multicolumn{2}{l}{\textbf{PRÄTERITUM} }\\
    ich & \textbf{ging}\\
    du & \textbf{gingst}\\
    er/sie/es & \textbf{ging}\\
    wir & \textbf{gingen}\\
    ihr & \textbf{gingt}\\
    sie/Sie & \textbf{gingen}\\
  \end{tabular}
  \hfill
  \begin{tabular}{rl}
     \multicolumn{2}{l}{\textbf{IMPERATIV} }\\
   & \textbf{geh(e) (du)!}\\
   & \textbf{gehen wir!}\\
   & \textbf{geht (ihr)!}\\
   & \textbf{gehen Sie!}\\
   & \vspace{10pt}
  \end{tabular}

    \vspace{10pt}

 \begin{tabular}{rl}
     \multicolumn{2}{l}{\textbf{PERFEKT}}\\
   & sein + \textbf{gegangen}\\
  \end{tabular}
\hfill
   \begin{tabular}{rl}
   
    \multicolumn{2}{l}{\textbf{FUTURE I} }\\
   & werden + \textbf{gehen}\\
 
  \end{tabular}
\hfill
   \begin{tabular}{rl}
      \multicolumn{2}{l}{\textbf{PLUSQUAMPERFEKT}}\\
   & waren + \textbf{gegangen}\\
  \end{tabular}
  
  
\vspace{10pt}

\begin{tabular}{lll}
\multicolumn{3}{l}{\textbf{MIT PRÄPOSITIONEN}}\\
& \textbf{---} & \\
\end{tabular}
\hfill
\begin{tabular}{lll}
\multicolumn{3}{l}{\textbf{TRENNBARE VERBEN}}\\
& \textbf{---} & \\
\end{tabular}
\vspace{-5pt}
\end{verbentry}

% gelten (to be valid / to apply)
\begin{verbentry}{
  style = {default},
  toc = {gelten (to be valid / to apply)},
  name = {gelten},
  english = {to be valid / to apply},
  type = {irregular},
  auxiliary = {haben}
}
 
    
     \vspace{5pt}

    \hrule
    
    \vspace{5pt}

  \begin{tabular}{rl}
     \multicolumn{2}{l}{\textbf{PRÄSENS} }\\
    ich & \textbf{gelte}\\
    du & \textbf{giltst}\\
    er/sie/es & \textbf{gilt}\\
    wir & \textbf{gelten}\\
    ihr & \textbf{geltet}\\
    sie/Sie & \textbf{gelten}\\
  \end{tabular}
  \hfill
     \begin{tabular}{rl}
    \multicolumn{2}{l}{\textbf{PRÄTERITUM} }\\
    ich & \textbf{galt}\\
    du & \textbf{galtest}\\
    er/sie/es & \textbf{galt}\\
    wir & \textbf{galten}\\
    ihr & \textbf{galtet}\\
    sie/Sie & \textbf{galten}\\
  \end{tabular}
  \hfill
  \begin{tabular}{rl}
     \multicolumn{2}{l}{\textbf{IMPERATIV} }\\
   & \textbf{gilt (du)!}\\
   & \textbf{gelten wir!}\\
   & \textbf{geltet (ihr)!}\\
   & \textbf{gelten Sie!}\\
   & \vspace{10pt}
  \end{tabular}

    \vspace{10pt}

 \begin{tabular}{rl}
     \multicolumn{2}{l}{\textbf{PERFEKT}}\\
   & haben + \textbf{gegolten}\\
  \end{tabular}
\hfill
   \begin{tabular}{rl}
    \multicolumn{2}{l}{\textbf{FUTURE I} }\\
   & werden + \textbf{gelten}\\
  \end{tabular}
\hfill
   \begin{tabular}{rl}
      \multicolumn{2}{l}{\textbf{PLUSQUAMPERFEKT}}\\
   & hatten + \textbf{gegolten}\\
  \end{tabular}

   \vspace{10pt}

   \begin{tabular}{lll}
      \multicolumn{3}{l}{\textbf{MIT PRÄPOSITIONEN}}\\
& \textbf{gelten für} + Akk & (to apply to / be valid for)\\
& \textbf{gelten als} + Nom & (to be considered as)\\
\end{tabular}
  \hfill
   \begin{tabular}{lll}
\multicolumn{3}{l}{\textbf{TRENNBARE VERBEN}}\\
& \textbf{---} & \\
\end{tabular}

  \vspace{-5pt}
\end{verbentry}

% gießen (to pour / to water)
\begin{verbentry}{
  style = {default},
  toc = {gießen (to pour / to water)},
  name = {gießen},
  english = {to pour / to water},
  type = {irregular},
  auxiliary = {haben}
}
 
    
     \vspace{5pt}

    \hrule
    
    \vspace{5pt}

  \begin{tabular}{rl}
     \multicolumn{2}{l}{\textbf{PRÄSENS} }\\
    ich & \textbf{gieße}\\
    du & \textbf{gießt}\\
    er/sie/es & \textbf{gießt}\\
    wir & \textbf{gießen}\\
    ihr & \textbf{gießt}\\
    sie/Sie & \textbf{gießen}\\
  \end{tabular}
  \hfill
     \begin{tabular}{rl}
    \multicolumn{2}{l}{\textbf{PRÄTERITUM} }\\
    ich & \textbf{goss}\\
    du & \textbf{gossest}\\
    er/sie/es & \textbf{goss}\\
    wir & \textbf{gossen}\\
    ihr & \textbf{gosst}\\
    sie/Sie & \textbf{gossen}\\
  \end{tabular}
  \hfill
  \begin{tabular}{rl}
     \multicolumn{2}{l}{\textbf{IMPERATIV} }\\
   & \textbf{gieß(e) (du)!}\\
   & \textbf{gießen wir!}\\
   & \textbf{gießt (ihr)!}\\
   & \textbf{gießen Sie!}\\
   & \vspace{10pt}
  \end{tabular}

    \vspace{10pt}

 \begin{tabular}{rl}
     \multicolumn{2}{l}{\textbf{PERFEKT}}\\
   & haben + \textbf{gegossen}\\
  \end{tabular}
\hfill
   \begin{tabular}{rl}
   
    \multicolumn{2}{l}{\textbf{FUTURE I} }\\
   & werden + \textbf{gießen}\\
 
  \end{tabular}
\hfill
   \begin{tabular}{rl}
      \multicolumn{2}{l}{\textbf{PLUSQUAMPERFEKT}}\\
   & hatten + \textbf{gegossen}\\
  \end{tabular}

   \vspace{10pt}

   \begin{tabular}{lll}
      \multicolumn{3}{l}{\textbf{MIT PRÄPOSITIONEN}}\\
& \textbf{gießen in} + Akk & (to pour into)\\
& \textbf{gießen über} + Akk & (to pour over)\\
\end{tabular}
  \hfill
   \begin{tabular}{lll}
      \multicolumn{3}{l}{\textbf{TRENNBARE VERBEN}}\\
 & \textbf{ein$|$gießen} & (to pour in)\\
 & \textbf{aus$|$gießen} & (to pour out)\\
\vspace{10pt}
  \end{tabular}


  \vspace{-5pt}
\end{verbentry}

% genießen (to enjoy)
\begin{verbentry}{
  style = {acc},
  toc = {genießen (to enjoy)},
  name = {genießen},
  english = {to enjoy},
  type = {irregular, + Akkusativ},
  auxiliary = {haben}
}
 
    
     \vspace{5pt}

    \hrule
    
    \vspace{5pt}

  \begin{tabular}{rl}
     \multicolumn{2}{l}{\textbf{PRÄSENS} }\\
    ich & \textbf{genieße}\\
    du & \textbf{genießt}\\
    er/sie/es & \textbf{genießt}\\
    wir & \textbf{genießen}\\
    ihr & \textbf{genießt}\\
    sie/Sie & \textbf{genießen}\\
  \end{tabular}
  \hfill
     \begin{tabular}{rl}
    \multicolumn{2}{l}{\textbf{PRÄTERITUM} }\\
    ich & \textbf{genoss}\\
    du & \textbf{genossest}\\
    er/sie/es & \textbf{genoss}\\
    wir & \textbf{genossen}\\
    ihr & \textbf{genosst}\\
    sie/Sie & \textbf{genossen}\\
  \end{tabular}
  \hfill
  \begin{tabular}{rl}
     \multicolumn{2}{l}{\textbf{IMPERATIV} }\\
   & \textbf{genieß(e) (du)!}\\
   & \textbf{genießen wir!}\\
   & \textbf{genießt (ihr)!}\\
   & \textbf{genießen Sie!}\\
   & \vspace{10pt}
  \end{tabular}

    \vspace{10pt}

 \begin{tabular}{rl}
     \multicolumn{2}{l}{\textbf{PERFEKT}}\\
  &  haben + \textbf{genossen}\\
  \end{tabular}
\hfill
   \begin{tabular}{rl}
   
    \multicolumn{2}{l}{\textbf{FUTURE I} }\\
  &  werden + \textbf{genießen}\\
 
  \end{tabular}
\hfill
   \begin{tabular}{rl}
      \multicolumn{2}{l}{\textbf{PLUSQUAMPERFEKT}}\\
  &  hatten + \textbf{genossen}\\
  \end{tabular}
  
  
\vspace{10pt}

\begin{tabular}{lll}
\multicolumn{3}{l}{\textbf{MIT PRÄPOSITIONEN}}\\
& \textbf{---} & \\
\end{tabular}
\hfill
\begin{tabular}{lll}
\multicolumn{3}{l}{\textbf{TRENNBARE VERBEN}}\\
& \textbf{---} & \\
\end{tabular}
\vspace{-5pt}
\end{verbentry}

% gestikulieren (to gesticulate)
\begin{verbentry}{
  style = {default},
  toc = {gestikulieren (to gesticulate)},
  name = {gestikulieren},
  english = {to gesticulate},
  type = {regular},
  auxiliary = {haben}
}


\vspace{5pt}
\hrule
\vspace{5pt}

\begin{tabular}{rl}
\multicolumn{2}{l}{\textbf{PRÄSENS}}\\
ich & \textbf{gestikuliere}\\
du & \textbf{gestikulierst}\\
er/sie/es & \textbf{gestikuliert}\\
wir & \textbf{gestikulieren}\\
ihr & \textbf{gestikuliert}\\
sie/Sie & \textbf{gestikulieren}\\
\end{tabular}
\hfill
\begin{tabular}{rl}
\multicolumn{2}{l}{\textbf{PRÄTERITUM}}\\
ich & \textbf{gestikulierte}\\
du & \textbf{gestikuliertest}\\
er/sie/es & \textbf{gestikulierte}\\
wir & \textbf{gestikulierten}\\
ihr & \textbf{gestikuliertet}\\
sie/Sie & \textbf{gestikulierten}\\
\end{tabular}
\hfill
\begin{tabular}{rl}
\multicolumn{2}{l}{\textbf{IMPERATIV}}\\
& \textbf{gestikuliere (du)!}\\
& \textbf{gestikulieren wir!}\\
& \textbf{gestikuliert (ihr)!}\\
& \textbf{gestikulieren Sie!}\\
\end{tabular}

\vspace{10pt}

\begin{tabular}{rl}
\multicolumn{2}{l}{\textbf{PERFEKT}}\\
& haben + \textbf{gestikuliert}\\
\end{tabular}
\hfill
\begin{tabular}{rl}
\multicolumn{2}{l}{\textbf{FUTURE I}}\\
& werden + \textbf{gestikulieren}\\
\end{tabular}
\hfill
\begin{tabular}{rl}
\multicolumn{2}{l}{\textbf{PLUSQUAMPERFEKT}}\\
& hatten + \textbf{gestikuliert}\\
\end{tabular}

\vspace{10pt}

\begin{tabular}{lll}
\multicolumn{3}{l}{\textbf{MIT PRÄPOSITIONEN}}\\
& \textbf{gestikulieren mit} + Dat & (to gesticulate with)\\
\end{tabular}
\hfill
\begin{tabular}{lll}
\multicolumn{3}{l}{\textbf{TRENNBARE VERBEN}}\\
& \textbf{---} & \\
\end{tabular}
\vspace{-5pt}
\end{verbentry}

% gewinnen (to win)
\begin{verbentry}{
  style = {acc},
  toc = {gewinnen (to win)},
  name = {gewinnen},
  english = {to win},
  type = {irregular, + Akkusativ},
  auxiliary = {haben}
}
 
    
     \vspace{5pt}

    \hrule
    
    \vspace{5pt}

  \begin{tabular}{rl}
     \multicolumn{2}{l}{\textbf{PRÄSENS} }\\
    ich & \textbf{gewinne}\\
    du & \textbf{gewinnst}\\
    er/sie/es & \textbf{gewinnt}\\
    wir & \textbf{gewinnen}\\
    ihr & \textbf{gewinnt}\\
    sie/Sie & \textbf{gewinnen}\\
  \end{tabular}
  \hfill
     \begin{tabular}{rl}
    \multicolumn{2}{l}{\textbf{PRÄTERITUM} }\\
    ich & \textbf{gewann}\\
    du & \textbf{gewannst}\\
    er/sie/es & \textbf{gewann}\\
    wir & \textbf{gewannen}\\
    ihr & \textbf{gewannt}\\
    sie/Sie & \textbf{gewannen}\\
  \end{tabular}
  \hfill
  \begin{tabular}{rl}
     \multicolumn{2}{l}{\textbf{IMPERATIV} }\\
   & \textbf{gewinn(e) (du)!}\\
   & \textbf{gewinnen wir!}\\
   & \textbf{gewinnt (ihr)!}\\
   & \textbf{gewinnen Sie!}\\
   & \vspace{10pt}
  \end{tabular}

    \vspace{10pt}

 \begin{tabular}{rl}
     \multicolumn{2}{l}{\textbf{PERFEKT}}\\
  &  haben + \textbf{gewonnen}\\
  \end{tabular}
\hfill
   \begin{tabular}{rl}
   
    \multicolumn{2}{l}{\textbf{FUTURE I} }\\
  &  werden + \textbf{gewinnen}\\
 
  \end{tabular}
\hfill
   \begin{tabular}{rl}
      \multicolumn{2}{l}{\textbf{PLUSQUAMPERFEKT}}\\
  &  hatten + \textbf{gewonnen}\\
  \end{tabular}
  
  
\vspace{10pt}

\begin{tabular}{lll}
\multicolumn{3}{l}{\textbf{MIT PRÄPOSITIONEN}}\\
& \textbf{---} & \\
\end{tabular}
\hfill
\begin{tabular}{lll}
\multicolumn{3}{l}{\textbf{TRENNBARE VERBEN}}\\
& \textbf{---} & \\
\end{tabular}
\vspace{-5pt}
\end{verbentry}

% glauben (to believe)
\begin{verbentry}{
  style = {dat},
  toc = {glauben (to believe)},
  name = {glauben},
  english = {to believe},
  type = {regular, + Dativ},
  auxiliary = {haben}
}
 
    
     \vspace{5pt}

    \hrule
    
    \vspace{5pt}

  \begin{tabular}{rl}
     \multicolumn{2}{l}{\textbf{PRÄSENS} }\\
    ich & \textbf{glaube}\\
    du & \textbf{glaubst}\\
    er/sie/es & \textbf{glaubt}\\
    wir & \textbf{glauben}\\
    ihr & \textbf{glaubt}\\
    sie/Sie & \textbf{glauben}\\
  \end{tabular}
  \hfill
     \begin{tabular}{rl}
    \multicolumn{2}{l}{\textbf{PRÄTERITUM} }\\
    ich & \textbf{glaubte}\\
    du & \textbf{glaubtest}\\
    er/sie/es & \textbf{glaubte}\\
    wir & \textbf{glaubten}\\
    ihr & \textbf{glaubtet}\\
    sie/Sie & \textbf{glaubten}\\
  \end{tabular}
  \hfill
  \begin{tabular}{rl}
     \multicolumn{2}{l}{\textbf{IMPERATIV} }\\
   & \textbf{glaub(e) (du)!}\\
   & \textbf{glauben wir!}\\
   & \textbf{glaubt (ihr)!}\\
   & \textbf{glauben Sie!}\\
   & \vspace{10pt}
  \end{tabular}

    \vspace{10pt}

 \begin{tabular}{rl}
     \multicolumn{2}{l}{\textbf{PERFEKT}}\\
   & haben + \textbf{geglaubt}\\
  \end{tabular}
\hfill
   \begin{tabular}{rl}
   
    \multicolumn{2}{l}{\textbf{FUTURE I} }\\
   & werden + \textbf{glauben}\\
 
  \end{tabular}
\hfill
   \begin{tabular}{rl}
      \multicolumn{2}{l}{\textbf{PLUSQUAMPERFEKT}}\\
   & hatten + \textbf{geglaubt}\\
  \end{tabular}
  
\vspace{10pt}

\begin{tabular}{lll}
\multicolumn{3}{l}{\textbf{MIT PRÄPOSITIONEN}}\\
& \textbf{glauben an} + Akk & (to believe in)\\
\end{tabular}
\hfill
\begin{tabular}{lll}
\multicolumn{3}{l}{\textbf{TRENNBARE VERBEN}}\\
& \textbf{---} & \\
\end{tabular}

  \vspace{-5pt}
\end{verbentry}

% grillen (to grill / barbecue)
\begin{verbentry}{
  style = {acc},
  toc = {grillen (to grill / barbecue)},
  name = {grillen},
  english = {to grill / barbecue},
  type = {regular, + Akkusativ},
  auxiliary = {haben}
}


\vspace{5pt}
\hrule
\vspace{5pt}

\begin{tabular}{rl}
\multicolumn{2}{l}{\textbf{PRÄSENS}}\\
ich & \textbf{grille}\\
du & \textbf{grillst}\\
er/sie/es & \textbf{grillt}\\
wir & \textbf{grillen}\\
ihr & \textbf{grillt}\\
sie/Sie & \textbf{grillen}\\
\end{tabular}
\hfill
\begin{tabular}{rl}
\multicolumn{2}{l}{\textbf{PRÄTERITUM}}\\
ich & \textbf{grillte}\\
du & \textbf{grilltest}\\
er/sie/es & \textbf{grillte}\\
wir & \textbf{grillten}\\
ihr & \textbf{grilltet}\\
sie/Sie & \textbf{grillten}\\
\end{tabular}
\hfill
\begin{tabular}{rl}
\multicolumn{2}{l}{\textbf{IMPERATIV}}\\
& \textbf{grille (du)!}\\
& \textbf{grillen wir!}\\
& \textbf{grillt (ihr)!}\\
& \textbf{grillen Sie!}\\
\end{tabular}

\vspace{10pt}

\begin{tabular}{rl}
\multicolumn{2}{l}{\textbf{PERFEKT}}\\
& haben + \textbf{gegrillt}\\
\end{tabular}
\hfill
\begin{tabular}{rl}
\multicolumn{2}{l}{\textbf{FUTURE I}}\\
& werden + \textbf{grillen}\\
\end{tabular}
\hfill
\begin{tabular}{rl}
\multicolumn{2}{l}{\textbf{PLUSQUAMPERFEKT}}\\
& hatten + \textbf{gegrillt}\\
\end{tabular}

\vspace{10pt}

\begin{tabular}{lll}
\multicolumn{3}{l}{\textbf{MIT PRÄPOSITIONEN}}\\
& \textbf{grillen auf} + Dat & (to grill on)\\
\end{tabular}
\hfill
\begin{tabular}{lll}
\multicolumn{3}{l}{\textbf{TRENNBARE VERBEN}}\\
& an$|$grillen & (to lightly grill / start grilling)\\
\end{tabular}

\vspace{-5pt}
\end{verbentry}

% gurgeln (to gargle)
\begin{verbentry}{
  style = {default},
  toc = {gurgeln (to gargle)},
  name = {gurgeln},
  english = {to gargle},
  type = {regular},
  auxiliary = {haben}
}


\vspace{5pt}
\hrule
\vspace{5pt}

\begin{tabular}{rl}
\multicolumn{2}{l}{\textbf{PRÄSENS}}\\
ich & \textbf{gurgle}\\
du & \textbf{gurgelst}\\
er/sie/es & \textbf{gurgelt}\\
wir & \textbf{gurgeln}\\
ihr & \textbf{gurgelt}\\
sie/Sie & \textbf{gurgeln}\\
\end{tabular}
\hfill
\begin{tabular}{rl}
\multicolumn{2}{l}{\textbf{PRÄTERITUM}}\\
ich & \textbf{gurgelte}\\
du & \textbf{gurgeltest}\\
er/sie/es & \textbf{gurgelte}\\
wir & \textbf{gurgelten}\\
ihr & \textbf{gurgeltet}\\
sie/Sie & \textbf{gurgelten}\\
\end{tabular}
\hfill
\begin{tabular}{rl}
\multicolumn{2}{l}{\textbf{IMPERATIV}}\\
& \textbf{gurgle (du)!}\\
& \textbf{gurgeln wir!}\\
& \textbf{gurgelt (ihr)!}\\
& \textbf{gurgeln Sie!}\\
\end{tabular}

\vspace{10pt}

\begin{tabular}{rl}
\multicolumn{2}{l}{\textbf{PERFEKT}}\\
& haben + \textbf{gegurgelt}\\
\end{tabular}
\hfill
\begin{tabular}{rl}
\multicolumn{2}{l}{\textbf{FUTURE I}}\\
& werden + \textbf{gurgeln}\\
\end{tabular}
\hfill
\begin{tabular}{rl}
\multicolumn{2}{l}{\textbf{PLUSQUAMPERFEKT}}\\
& hatten + \textbf{gegurgelt}\\
\end{tabular}

\vspace{10pt}

\begin{tabular}{lll}
\multicolumn{3}{l}{\textbf{MIT PRÄPOSITIONEN}}\\
& \textbf{gurgeln mit} + Dat & (to gargle with)\\
\end{tabular}
\hfill
\begin{tabular}{lll}
\multicolumn{3}{l}{\textbf{TRENNBARE VERBEN}}\\
& \textbf{---} & \\
\end{tabular}
\vspace{-5pt}
\end{verbentry}
